% !PRINTED 373
\weberpage{373}
%<<CONT
one first expresses $x_1$ in terms of $x_3$ and substitutes this into the expression of $x$ in terms of $x_1$.

By (8), a composite equivalence is a \gesperrt{proper} or an \gesperrt{improper} one according as an even or an odd number of improper ones is found among the components.

If one composes the two substitutions
% [mirrored 2026-08-07] the German restored the equation number (12) the page prints here
\weq{12}{\begin{pmatrix}\alpha, & \beta\\ \gamma, & \delta\end{pmatrix}, \quad
\begin{pmatrix}\varepsilon\delta, & -\varepsilon\beta\\ -\varepsilon\gamma, & \varepsilon\alpha\end{pmatrix}}
then one obtains, no matter which of the two one puts in the first place,
\weq{13}{\begin{pmatrix}1, & 0\\ 0, & 1\end{pmatrix}.}

This substitution has, by (1), the same meaning as $y = x$; it changes nothing, and is called the \gesperrt{identical substitution}, and indeed is also denoted briefly by (1). Accordingly one calls the two substitutions (12) \gesperrt{reciprocal} to one another, and denotes them by
\[
S, \; S^{-1},
\]
or one sets
\weq{14}{\begin{pmatrix}\varepsilon\delta, & -\varepsilon\beta\\ -\varepsilon\gamma, & \varepsilon\alpha\end{pmatrix} = \begin{pmatrix}\alpha, & \beta\\ \gamma, & \delta\end{pmatrix}^{-1}.}

Under composition with the identical substitution (13) every other substitution remains unchanged.

From this we very simply derive the proof of the theorem:

\gesperrt{V. All rational numbers are equivalent to one another, and indeed both properly and improperly.}

For if $m:n$ and $m':n'$ are two rational fractions, and $m$ and $n$, as well as $m'$ and $n'$, are without common divisor, then, if $\varepsilon, \varepsilon'$ are $\pm 1$ at pleasure, we can determine the numbers $\mu, \nu$ and $\mu', \nu'$ according to \S.\,118 so that
\[
m\nu - n\mu = \varepsilon, \qquad m'\nu' - n'\mu' = \varepsilon'
\]
holds, and so that therefore
\[
M = \begin{pmatrix}m, & \mu\\ n, & \nu\end{pmatrix}, \quad M' = \begin{pmatrix}m', & \mu'\\ n', & \nu'\end{pmatrix}
\]
are two linear substitutions. It is then also
% [mirrored 2026-08-07] the German restored the dropped final display of page 373
\[
M M'^{-1} = \begin{pmatrix}\alpha, & \beta\\ \gamma, & \delta\end{pmatrix}
\]
%CONT>>
